% SPDX-License-Identifier: GPL-3.0-only
% Licensed under GNU GPL version 3; https://www.gnu.org/licenses/gpl-3.0.html . See LICENSE at repository root.
\section{Scope and conventions}\label{sec:overview}

LPL combines observations of a finite trace with modalities that extend that
trace by program executions. A formula can inspect its first state, move to
a proper suffix, or constrain every extension produced by a regular program.
A sequent assigns formulas to symbolic trace labels. A finite proof graph
may contain backlinks, provided the checked trace information supports a
well-founded descent argument.

The source development fixes four conventions that affect almost every rule.
First, label juxtaposition is ordinary list append, including repeated states.
Second, quantifier bodies belong to pure first-order logic. Third, one
quantified value is written uniformly throughout a trace. Fourth, sequent
truth is guarded by nonemptiness of every complete assertion label.
These conventions are built into the definitions below rather than supplied
as optional assumptions.

\paragraph{Relationship to the manuscript and Iris.}
The original mathematical presentation is in
\nolinkurl{Theory/20260926-Types-by-AI.tex}. The current Rocq source additionally supports
many-sorted signatures, reusable parallel updates, and nine primitive formula
constructors. This manual follows that current source. The presentation style
and modular file layout follow \code{iris/tex}; LPL has its own notation package
and does not depend on Iris's logic or LaTeX macros.

\paragraph{Notation.}
Letters $x,y,z$ denote program variables; $i$ denotes a world name; $k$ denotes
a trace name; $t,u$ denote terms; $A,B$ denote formulas; and $\alpha,\beta$ denote
programs. We use $r,q$ for reusable updates and $U,V$ for rooted state
expressions. Labels are $\ell,\sigma,\tau$. Program variables, world names,
trace names, tags, and graph-node indices occupy separate namespaces even
when all but the program variables are represented by natural numbers.

For lists, $\epsilon$ or $[]$ denotes the empty list, $s::\pi$ denotes cons,
$\cat$ denotes ordinary append, and $|\pi|$ denotes length. An execution
composition $\pi\odot\eta\leadsto\theta$ is a relation, defined in
Section~\ref{sec:traces}. Label juxtaposition $\jlabel$ denotes a syntax
constructor whose interpretation is append. These symbols will not be used
interchangeably.

\paragraph{What a theorem reference means.}
A statement headed ``Theorem'', ``Lemma'', or ``Proposition'' is a
mathematical restatement of the linked source results, with their relevant
premises made explicit. Proof sketches explain those results rather than
replace their Rocq proofs. Examples that merely unfold a definition are
identified as calculations. The manual itself is LaTeX, so maintaining a
source link does not mechanically prove that its prose is equivalent to Rocq.
The accompanying check records source hashes and validates declaration names.

\paragraph{Limits of the reference.}
The implementation proves soundness for the supported local instances and
accepted certificates. It supplies finite local rule generation and explicit
graph-building commands. It does not establish completeness, automatic proof
search, automatic first-order reasoning, or acceptance of every manuscript
preproof. Signature symbols do not acquire algebraic laws simply by being
declared; a model supplies their meanings and proofs of any desired laws.
